% =====================================================================
%  U02：PDF p.23--p.44（原书印刷 p.9--p.30）
%  1.2 Langlands parameters over the complex field（1.2.1--1.2.6）
%  1.3 Representations of loop groups（1.3.1--1.3.6）
%  本片段从节中部开始，不写 \chapter；节编号由 main.tex 中的第 1 章接续。
% =====================================================================

% 接续 U01 末段：p.23 第一句（“Some recent results ...” 的收尾）已由 U01 补译，
% 此处从 p.23 第二句起译；该句收尾已并入 ch01-1.tex 末段，此处不再重复。

\section{Langlands 参数（复域）}

\begin{keypoints}
    \item \textbf{核心思想：} 复域上的局部 Langlands 参数不取同态 $\hat\Z\to{}^L G$（这样的参数太少），而取去心圆盘上的平坦 ${}^L G$-丛，即联络 $\partial_t+A(t)$ 的规范等价类。
    \item \textbf{数学工具：} Galois 群与基本群（射有限完备化）；平坦丛与单值性表示；Tannakian 形式；全纯联络与 Riemann--Hilbert 对应。
    \item \textbf{关键结论：} 只允许正则奇点时，联络的等价类恰是 ${}^L G$ 的共轭类，与非正则奇点相比参数太少；允许任意奇点后得到的 $\Loc_{{}^L G}(D^\times)$ 才是合适的参数空间。
    \item \textbf{记号提醒：} $D^\times=\Spec\C((t))$；${}^L\fg=\Lie({}^L G)$；$A_{-1}$ 是 $A(t)$ 在 $t=0$ 处的留数；${}^L G^\circ$ 从 p.10 起直接简记为 ${}^L G$。
\end{keypoints}

% [Tutor Note: p.9 引言。把 $\F_q((t))$ 换成 $\C((t))$ 的动机是：复域上可以使用代数几何的工具，而局部 Langlands 对应的模式（表示论一侧对应参数一侧）在两种域上应当形式相同。注意这里只说“希望”，并未断言对应已经存在，真正的表述要到 p.30 才给出。]
现在我们要找局部 Langlands 猜想的一个推广：把域 $F=\F_q((t))$ 换成域 $\C((t))$。我们想看看 Langlands 对应的思想与模式在这个新背景下如何展开，希望借此更好地理解这一对应背后的深层结构。

从今以后，$G$ 都是 $\C$ 上的连通约化群，$G(F)$ 指群 $G((t))=G(\C((t)))$，也称为圈群；更精确地说，是形式圈群。这里“形式”一词指形式洛朗幂级数代数 $\C((t))$，与之相对的是群 $G(\C[t,t^{-1}])$，其中 $\C[t,t^{-1}]$ 是洛朗多项式代数。后一个群可以看成从单位圆 $|t|=1$ 到 $G$ 的映射，即 $G$ 中的“圈”。

于是，我们希望研究圈群 $G((t))$ 的光滑表示，并试图把它们与某些“Langlands 参数”联系起来。类比上面讨论的局部非阿基米德域的情形，我们期待这些参数与 $\C((t))$ 的 Galois 群以及 Langlands 对偶群 ${}^L G$ 有关。

本书讨论的圈群局部 Langlands 对应，可以看作把 Langlands 纲领的思想搬进复代数几何的第一步。特别地，它对整体几何 Langlands 对应有深远影响（详见第 10.5 节与 [Frenkel 2007]）。这其实也是本项目的一个动机。

\subsection{Galois 群与基本群}

% [Tutor Note: p.9--10 的结构。本节先说明“天真”的 Langlands 参数是什么，再论证它为什么不可用（同态 $\hat\Z\to{}^L G$ 必过有限商），最后用基本群与平坦丛的关系给出替代方案。全书的参数定义只依赖后一步。]
先描述 $F=\C((t))$ 的 Galois 群 $\operatorname{Gal}(\overline F/F)$。注意 $F$ 的代数闭包 $\overline F$ 同构于域 $\C((t^{1/n}))$（$n\ge 0$）沿自然包含 $\C((t^{1/n}))\hookrightarrow\C((t^{1/m}))$（当 $n\mid m$）的归纳极限。因此 $\operatorname{Gal}(\overline F/F)$ 是 Galois 群
\[
    \operatorname{Gal}(\C((t^{1/n}))/\C((t)))\simeq \Z/n\Z
\]
的反向极限，其中 $\overline{k}\in\Z/n\Z$ 对应 $\C((t^{1/n}))$ 上把 $t^{1/n}$ 送到 $e^{2\pi i k/n}t^{1/n}$ 的自同构。结论是
\[
    \operatorname{Gal}(\overline F/F)\simeq \hat\Z,
\]
其中 $\hat\Z$ 是 $\Z$ 的射有限完备化，我们已在 1.1.3 节遇到过。

不过要注意：在研究 $\F_q((t))$ 的 Galois 群时，$\hat\Z$ 是作为商群出现的，对应系数域 $\F_q$ 的 Galois 群。现在系数域是 $\C$，它代数闭，于是 $\hat\Z$ 就是 $\C((t))$ 的整个 Galois 群。

Langlands 参数的天真类比是同态 $\operatorname{Gal}(\overline F/F)\to{}^L G$ 的等价类，即同态 $\hat\Z\to{}^L G$。$G$ 定义在 $\C$ 上，因而它的所有极大环面都分裂，于是群 $G((t))$ 也含一个分裂环面 $T((t))$，其中 $T$ 是 $G$ 的极大环面（但它也含非分裂极大环面，因为域 $\C((t))$ 不是代数闭的）。因此 Langlands 对偶群 ${}^L G$ 是 Galois 群与群 ${}^L G^\circ$ 的直积。既然是直积，我们就可以——而且将会——把注意力限制在 ${}^L G^\circ$ 上。为简化记号，从今以后把 ${}^L G^\circ$ 直接记作 ${}^L G$。

同态 $\hat\Z\to{}^L G$ 必然经过一个有限商 $\hat\Z\to\Z/n\Z$ 分解。所以 $\hat\Z\to{}^L G$ 的同态等价类，与 ${}^L G$ 中有限阶元素的共轭类是同一回事。这样的元素太少，支撑不起 Langlands 对应的有意义的推广。因此我们寻找更合理的替代。

% [Tutor Note: p.10。这一步“太少”的具体机制是：$\hat\Z$ 射有限，任何到复代数群 ${}^L G$ 的连续同态都落在某个有限商 $\Z/n\Z$ 里，于是像集是有限阶元。共轭类的离散集合无法承载表示论中的连续族（例如非正则奇点带来的连续模），这正是后文引入联络的动机。]
回顾 Galois 群与基本群的联系。设 $X$ 是 $\C$ 上的代数簇。若 $Y\to X$ 是 $X$ 的覆叠，则 $Y$ 上有理函数域 $\C(Y)$ 是 $X$ 上有理函数域 $F=\C(X)$ 的扩域。覆叠的覆叠变换，即诱导 $X$ 上恒等映射的 $Y$ 的自同构，给出保持 $\C(X)\subset\C(Y)$ 的 $\C(Y)$ 自同构。于是把 Galois 群 $\operatorname{Gal}(\C(Y)/\C(X))$ 与覆叠变换群等同起来。若覆叠无分歧，则这个群可以等同于 $X$ 基本群的一个商；否则，它同构于把分歧除子挖掉后 $X$ 的基本群的商。

特别地，$\C(X)$ 的最大无分歧扩域（可以看成 $X$ 的“最大无分歧覆叠”的函数域）的 Galois 群，是 $X$ 的基本群 $\pi_1(X)$ 的射有限完备化。同样，对任意除子 $D\subset X$，$\C(X)$ 在 $D$ 之外无分歧的最大扩域的 Galois 群，是 $\pi_1(X\setminus D)$ 的射有限完备化，记作 $\pi_1^{\mathrm{alg}}(X\setminus D)$。$\C(X)$ 的代数闭包，是 $X$ 在各除子 $D\subset X$ 处分歧的最大覆叠的函数域沿除子包含所对应的自然包含取的归纳极限。因此 Galois 群 $\operatorname{Gal}(\overline{\C(X)}/\C(X))$ 是群 $\pi_1^{\mathrm{alg}}(X\setminus D)$ 沿映射 $\pi_1^{\mathrm{alg}}(X\setminus D')\to\pi_1^{\mathrm{alg}}(X\setminus D)$（$D\subset D'$）的反向极限。

严格说来，定义 $X$ 的基本群需要选一个基点 $x\in X$。上面的讨论中我们暗中取 $X$ 的泛点，即 $\Spec\C(X)\hookrightarrow X$，作为基点。但不同的 $x$ 给出的基本群彼此同构，因此它们表示的等价类（这才是我们关心的）也相同。所以我们忽略基点的选取。

% [Tutor Note: p.11。取泛点作基点是本书的简化。若不取基点，单值性表示只能在共轭意义下确定；对 $GL_n$ 这无碍（共轭类正是要的参数），但后面讲 ${}^L G$-丛时，“规范等价”与“共轭”的差别会显式出现，见 p.18 与式 (1.2-7)。]
回到我们的设定，可见 $F=\C((t))$ 的 $\operatorname{Gal}(\overline F/F)$ 正是 $\hat\Z$，即去心圆盘 $D^\times$ 的“拓扑”基本群的射有限完备化\transnote{$\C((t))$ 的绝对 Galois 群是 $\hat\Z$，而不是 $\Z$：代数闭包 $\overline{\C((t))}=\bigcup_{n\ge1}\C((t^{1/n}))$ 由全部 Kummer 覆叠拼接而成，故其 Galois 群是有限群 $\Z/n\Z$ 的反向极限。去心圆盘的拓扑基本群是 $\pi_1(D^\times)=\Z$，其射有限完备化同样是 $\hat\Z$，这正是“代数基本群等于拓扑基本群的射有限完备化”这一一般事实在本例中的体现。[来源：原书 p.9--11]}。于是，天真的 Langlands 参数对应同态 $\pi_1^{\mathrm{alg}}(D^\times)\to{}^L G$。但在复的情形，同态 $\pi_1(X)\to{}^L G$ 可以从几何上更精细的数据得到，即从带平坦联络的丛得到；下面就来讨论。

\subsection{平坦丛}

设 $X$ 是实流形，$E$ 是 $X$ 上的复向量丛。于是给定映射 $p:E\to X$，满足条件：每点 $x\in X$ 有邻域 $U$，使得 $U$ 在 $p$ 下的原像同构于 $U\times\C^n$。但 $E$ 在 $U$ 上的这种平凡化不唯一；两个不同平凡化相差一个光滑映射 $U\to GL_n(\C)$。更具体地描述 $E$：取 $X$ 的一族开覆盖 $U_\alpha$，$\alpha\in A$，并在每个 $U_\alpha$ 上取 $E$ 的平凡化。则两个诱导平凡化在交叠 $U_\alpha\cap U_\beta$ 上的差别由一个函数 $g_{\alpha\beta}:U_\alpha\cap U_\beta\to GL_n(\C)$ 给出，称为转移函数。转移函数满足重要的传递性：在三重交叠 $U_\alpha\cap U_\beta\cap U_\gamma$ 上有 $g_{\alpha\gamma}=g_{\alpha\beta}g_{\beta\gamma}$。

这描述的是 $X$ 上普通向量丛的数据。平坦向量丛 $E$ 的数据，则是 $E$ 在任意点的足够小开邻域上的一族“优先”平凡化，它们构成传递系统，并且两个优先平凡化之差是取常值的 $GL_n(\C)$ 函数。具体说，平坦丛可以用 $X$ 的一族足够细的开覆盖 $U_\alpha$ 来描述：在每个 $U_\alpha$ 上取属于这个优先系统的平凡化。此时转移函数 $g_{\alpha\beta}:U_\alpha\cap U_\beta\to GL_n(\C)$ 是常函数，并满足传递性。这与普通丛的情形形成对照：普通丛的转移函数可以是任意的光滑 $GL_n(\C)$ 函数。

从微分几何的角度看，优先平凡化系统的数据可以简洁地用平坦联络的数据来表达。后者正是对 $E$ 的截面作微分所需的数据。

更精确地说，设 $\cO_X$ 是 $X$ 上光滑函数层，$\cT$ 是 $X$ 上光滑向量场层。用 $\End(E)$ 记 $E$ 的自同态丛；为简便，也用同一记号表示其光滑截面层。$E$ 上的一个联络是映射
\[
    \nabla:\cT\longrightarrow\End(E),
\]
它把开集 $U\subset X$ 上定义的向量场 $\xi\in\cT(U)$ 送到 $E$ 在 $U$ 上的自同态 $\nabla_\xi$。它应具有以下性质：
\begin{itemize}
    \item 它是 $\cO_X$ 线性的：$\nabla_{\xi+\eta}=\nabla_\xi+\nabla_\eta$，且 $\nabla_{f\xi}=f\nabla_\xi$，其中 $f\in\cO_X$。
    \item 它满足 Leibniz 法则：$\nabla_\xi(f\phi)=f\nabla_\xi(\phi)+\phi\cdot\xi(f)$，其中 $f\in\cO_X$，$\phi\in\End(E)$。
\end{itemize}

% [Tutor Note: p.12 公式 (1.2-1)。这里的 $\nabla$ 取值在 $\End(E)$，而 $\End(E)$ 带自然的李代数结构（交换子），$\cT$ 也带向量场的李代数结构；(1.2-1) 说的就是“联络是李代数同态”，等价于曲率为零。符号约定：$\nabla_{\xi+\eta}=\nabla_\xi+\nabla_\eta$ 与 Leibniz 法则合起来表明 $\nabla$ 是 $\cO_X$-线性的一阶微分算子。]
若联络还满足：$\nabla$ 是李代数同态
\begin{equation}
    [\nabla_\xi,\nabla_\chi]=\nabla_{[\xi,\chi]},
    \tag{1.2-1}
\end{equation}
其中 $\cT$ 与 $\End(E)$ 取自然的李代数结构，则称该联络是平坦的。

下面讨论联络的具体实现。在足够小的开集 $U$ 上可取坐标 $x_i$，$i=1,\dots,N$，并把丛 $E$ 平凡化。由 $\cO_X$ 线性，要确定联络 $\nabla$ 在 $U$ 上的限制，只需写出算子 $\nabla_{\partial_{x_i}}$，这里 $\partial_{x_i}=\partial/\partial x_i$。Leibniz 法则表明这些算子必有形式
\[
    \nabla_{\partial_{x_i}}=\partial_{x_i}+A_i(\mathbf{x}),
\]
其中 $A_i$ 是 $U$ 上 $n\times n$ 矩阵值光滑函数。

平坦性条件于是写成
\begin{equation}
    [\nabla_{\partial_{x_i}},\nabla_{\partial_{x_j}}]
    =\partial_{x_i}A_j-\partial_{x_j}A_i+[A_i,A_j]=0.
    \tag{1.2-2}
\end{equation}

所以，在实的情形要具体描述一个平坦联络，只需局部给出满足上述方程的矩阵 $A_i$。但还须保证它们在坐标变换与丛 $E$ 的平凡化变换下按正确的方式变换\transnote{以下两式原文各有一处笔误，译文照原书保留。其一，本节以 $N$ 记坐标个数、以 $n$ 记 $E$ 的秩，两处求和的上限应作 $N$ 而非 $n$。其二，第二式左端应作 $\nabla_{\partial_{y_i}}$：自由指标须与右端的 $i$ 一致，否则等式两边指标不匹配。[来源：原书 p.12 图版]}。若引入新坐标使 $x_i=x_i(\mathbf{y})$，则
\[
    \partial_{y_i}=\sum_{j=1}^{n}\frac{\partial x_j}{\partial y_i}\partial_{x_j},
\]
从而
\[
    \nabla_{\partial_{y_j}}=\sum_{j=1}^{n}\frac{\partial x_j}{\partial y_i}\nabla_{\partial_{x_j}}.
\]

% [Tutor Note: p.12--13 两处笔误，译文照原书保留。第一，指标个数：本节用 $N$ 表示坐标个数、$n$ 表示 $E$ 的秩，但这里两个求和都写成 $\sum_{j=1}^{n}$，应作 $\sum_{j=1}^{N}$。第二，第二个式子左端写成 $\nabla_{\partial_{y_j}}$，而右端自由指标是 $i$，应作 $\nabla_{\partial_{y_i}}$。此处极易误抄，请看原文图像核对。]
% [Tutor Note: p.13 规范变换的符号约定。$g\nabla_{\partial_{x_i}}g^{-1}=\nabla_{\partial_{x_i}}+gA_ig^{-1}-(\partial_{x_i}g)g^{-1}$ 来自把平凡化换成 $g^{-1}$ 后对 $\nabla(g^{-1}f)$ 用 Leibniz 法则；注意 $(n\times n)$ 矩阵的乘法次序不能交换，最后一项是 $-(\partial g)g^{-1}$ 而非 $g^{-1}\partial g$。这是全章符号错误的常见来源，后文 (1.2-4) 采用同一约定。]
给定 $E$ 的一个截面，它相对于旧平凡化表现为 $U$ 上的 $\C^n$ 值函数 $f(\mathbf{x})$；相对于新平凡化则表现为函数 $g(\mathbf{x})\cdot f(\mathbf{x})$，其中 $g(\mathbf{x})$ 是 $GL_n(\C)$ 值的转移函数。因此联络算子变成
\[
    g\nabla_{\partial_{x_i}}g^{-1}
    =\nabla_{\partial_{x_i}}+gA_ig^{-1}-(\partial_{x_i}g)g^{-1}.
\]

利用联络算子，可以如下构造丛的一族传递的优先局部平凡化：局部地，在小开集 $U$ 上考虑一阶微分方程组
\begin{equation}
    \nabla_\xi\Phi(\mathbf{x})=0,\qquad \xi\in\cT(U).
    \tag{1.2-3}
\end{equation}
这个方程组要有意义，就须满足平坦性条件 (1.2-2)。（否则 (1.2-2) 的右端是一个矩阵值函数，它必须零化 $\Phi(\mathbf{x})$；这样就得把 $\Phi(\mathbf{x})$ 限制到这个矩阵值函数的核上。）在满足条件时，线性微分方程解的存在唯一性定理告诉我们：对每一点 $\mathbf{x}_0\in U$ 和纤维 $E_{\mathbf{x}_0}$ 中的每个向量 $v$，方程组有唯一的局部解 $\Phi_v$ 满足初值条件 $\Phi(\mathbf{x}_0)=v$。不同的 $v\in E_{\mathbf{x}_0}$ 给出的解 $\Phi_v$ 是 $E$ 在 $U$ 上的截面，称为水平的。它们给出 $E$ 在 $U$ 各点处纤维的一族传递的优先等同：$E_{\mathbf{x}}\simeq E_{\mathbf{y}}$，$\mathbf{x},\mathbf{y}\in U$。若把其中一个纤维（例如 $\mathbf{x}_0$ 处的纤维）与 $\C^n$ 等同，就得到 $E$ 在 $U$ 上的一个平凡化。改变等同 $E_{\mathbf{x}_0}\simeq\C^n$，只会把这个平凡化改变一个常函数 $U\to GL_n$。于是得到 $E$ 所需的优先平凡化系统。

换一种说法：平坦联络的数据给出邻近纤维的一族优先等同，它们在局部是传递的。但若用它们沿 $X$ 中从 $\mathbf{x}_0$ 出发又回到 $\mathbf{x}_0$ 的给定道路来等同 $E$ 的纤维，就可能得到纤维 $E_{\mathbf{x}_0}$ 的一个非平凡自同构，即平坦联络沿这条道路的单值性。（更具体地说，这个单值性可以通过考虑方程组 (1.2-3) 在道路的小管状邻域内的解来求得。）这个自同构只依赖于道路的同伦类，所以 $E$ 上任何平坦联络都给出同态 $\pi_1(X,\mathbf{x}_0)\to\Aut E_{\mathbf{x}_0}$，其中 $\Aut E_{\mathbf{x}_0}$ 是 $E_{\mathbf{x}_0}$ 的自同构群。若把 $E_{\mathbf{x}_0}$ 与 $\C^n$ 等同，就得到同态 $\pi_1(X,\mathbf{x}_0)\to GL_n(\C)$。

% [Tutor Note: p.13。单值性给出的不只是一个同态，而是同态的共轭类：换基点 $\mathbf{x}_0$ 或换 $E_{\mathbf{x}_0}\simeq\C^n$ 的等同，只会把表示共轭掉。反过来，同态的共轭类由水平截面系统唯一决定。因此“平坦丛的同构类”与“$\pi_1$ 表示的共轭类”之间的对应是双向的（对 $GL_n$ 即 Riemann--Hilbert 对应），这是后文把 Langlands 参数解释为平坦丛的依据。]
于是，我们把 $X$ 上秩为 $n$ 的平坦丛对应到同态 $\pi_1(X,\mathbf{x}_0)\to GL_n(\C)$ 的一个等价类。

当 $X$ 是去心圆盘 $D^\times$，或光滑射影复曲线时，同态 $\pi_1(X,\mathbf{x}_0)\to GL_n(\C)$ 的等价类正是我们在复的情形下要当作 Langlands 参数的候选对象。这就是我们关心它们以带平坦联络的丛的形式重新表述的原因。

\subsection{全纯情形下的平坦丛}

% [Tutor Note: p.14。全纯丛的“全纯结构”由 $\bar\partial$-算子 $\nabla_{\partial_{\bar z_i}}$ 决定；平坦联络还缺 $\nabla_{\partial_{z_i}}$ 这一半。若 $X$ 是曲线（$N=1$），任何全纯联络都自动平坦：曲率是 2-形式，而曲线上 2-形式空间为零。这是维数上的事实，不涉及丛的具体形式。]
现在设 $X$ 是复代数簇，$E$ 是 $X$ 上的全纯向量丛。回忆：给 $X$ 上的向量丛 $E$ 一个全纯向量丛结构，等价于指定 $E$ 的哪些截面算作全纯的。在 $X$ 的开集 $U$ 上取全纯坐标 $z_i$，$i=1,\dots,N$，以及反全纯坐标 $\overline{z}_i$，$i=1,\dots,N$；指定全纯截面等价于在 $E$ 的截面上定义 $\bar\partial$-算子 $\nabla_{\partial_{\overline{z}_i}}$，$i=1,\dots,N$。于是 $E$ 的全纯截面恰好是这些算子零化的截面。这样我们就得到反全纯向量场在 $E$ 的一切光滑截面上的作用。这个作用给出平坦联络数据的“一半”。所以，要在全纯向量丛 $E$ 上指定平坦联络，只需在 $E$ 的全纯截面上定义 $X$ 的（局部）全纯向量场的作用。

在 $X$ 的足够小开集 $U$ 上，现在可以取与全纯结构相容的平凡化，使 $\nabla_{\partial_{\overline{z}_i}}=\partial_{\overline{z}_i}$。把这些数据扩充为 $E$ 上平坦联络的数据，就归结为构造算子
\[
    \nabla_{\partial_{z_i}}=\partial_{z_i}+A_i(\mathbf{z}),
\]
其中平坦性条件要求 $A_i$ 是 $z_j$（$j=1,\dots,N$）的全纯矩阵值函数，且
\[
    [\nabla_{\partial_{z_i}},\nabla_{\partial_{z_j}}]
    =\partial_{z_i}A_j-\partial_{z_j}A_i+[A_i,A_j]=0.
\]

特别地，若 $X$ 是复曲线，即 $N=1$，则全纯向量丛上的任何全纯联络都给出平坦联络，从而给出同态 $\pi_1(X)\to GL_n(\C)$ 的一个等价类。\footnote{不过要注意，曲线上有些全纯向量丛根本不带任何全纯联络；例如，线丛带全纯联络当且仅当其次数为 $0$。}

关于向量丛（等价地，主 $GL_n$-丛）上平坦联络的讨论到此结束。下面定义主 $G$-丛上的平坦联络，其中 $G$ 是复代数群。

\subsection{平坦 \texorpdfstring{$G$}{G}-丛}

回忆：流形 $X$ 上的主 $G$-丛是一个纤维化于 $X$ 之上的流形 $\mathcal{P}$，使得 $G$ 在 $\mathcal{P}$ 上有自然的纤维保持右作用，且在每个纤维上单可递。此外它是局部平凡的：每点有足够小的开邻域 $U\subset X$，$\mathcal{P}$ 在其上可平凡化，即存在与 $G$ 的右作用交换的同构 $t_U:\mathcal{P}|_U\xrightarrow{\ \sim\ }U\times G$。更具体地说，可取 $X$ 的开覆盖 $U_\alpha$，$\alpha\in A$，并在每个 $U_\alpha$ 上取 $\mathcal{P}$ 的平凡化。则两个诱导平凡化在交叠 $U_\alpha\cap U_\beta$ 上的差别由转移函数 $g_{\alpha\beta}:U_\alpha\cap U_\beta\to G$ 描述，它们满足传递性条件，与 $GL_n$ 的情形一样。

这样的丛 $\mathcal{P}$ 上的平坦结构，是 $\mathcal{P}$ 的各纤维之间一族局部传递的优先等同，等价地，是一族传递的优先局部平凡化（即转移函数 $g_{\alpha\beta}$ 是取常值的 $G$ 值函数）。为了给出平坦联络的微分几何实现，我们使用 Tannakian 形式。粗略地说，它断言：$G$ 可由其有限维表示范畴（带张量积结构并满足各种相容性）以及到向量空间范畴的“纤维函子”重构出来。

% [Tutor Note: p.15。Tannakian 重构需要“纤维函子”到向量空间范畴；在本节情形，这个函子就是 $V\mapsto V_{\mathcal{P}}$。这里只用到“张量函子可重构主丛”这一方向，省略了 Deligne 的一般定理与各种相容性（结合约束、交换约束）的验证，可参看 [Deligne and Milne 1982] 或 [Milne 2017]。]
给定 $X$ 上的主 $G$-丛 $\mathcal{P}$ 与 $G$ 的有限维表示 $V$，可以构造相伴向量丛
\[
    V_{\mathcal{P}}=\mathcal{P}\times_G V.
\]
于是得到从 $G$ 的有限维表示范畴到 $X$ 上向量丛范畴的函子。这两个范畴都是张量范畴，且上述函子显然是张量函子。Tannakian 形式使我们能从这个函子重构主 $G$-丛 $\mathcal{P}$。换句话说，$X$ 上主 $G$-丛的数据，等价于 $X$ 上一族以 $G$ 的有限维表示为标号的向量丛，连同对每对表示 $V,W$ 的同构 $(V\otimes W)_{\mathcal{P}}\simeq V_{\mathcal{P}}\otimes W_{\mathcal{P}}$，以及由该函子的张量性质带来的各种相容性。

现在，要在全纯主 $G$-丛 $\mathcal{P}$ 上定义平坦联络，就需要以相容的方式在每个相伴向量丛 $V_{\mathcal{P}}$ 上定义平坦联络。这意味着：在每个足够小的开集 $U\subset X$ 上，取 $U$ 上的一族局部全纯坐标 $z_i$，$i=1,\dots,N$，并取 $\mathcal{P}$ 在 $U$ 上的全纯平凡化之后，$\mathcal{P}$ 上平坦联络的数据由微分算子
\[
    \nabla_{\partial_{z_i}}=\partial_{z_i}+A_i(\mathbf{z})
\]
给出，其中 $A_i$ 现在是 $U$ 上取值于 $G$ 的李代数 $\fg$ 的全纯函数。

算子 $\nabla_{\partial_{z_i}}$ 在坐标变换和平凡化变换下的变换性质，由与上面相同的公式给出。特别地，在由 $U$ 上全纯 $G$ 值函数 $g$ 给出的平凡化变换下，
\begin{equation}
    \partial_{z_i}+A_i\longmapsto \partial_{z_i}+gA_ig^{-1}-(\partial_{z_i}g)g^{-1}.
    \tag{1.2-4}
\end{equation}
这样的变换称为规范变换。

% [Tutor Note: p.16。公式 (1.2-4) 的右端逐项都在 $\End V$ 中有定义，对 $G$ 的每个有限维表示 $V$ 都成立；Tannakian 形式保证存在 $\fg$ 中的“万有”元素，使其在任意 $V$ 上的作用恰为 $gA_ig^{-1}-(\partial_{z_i}g)g^{-1}$。$\fg$ 值是靠“对所有 $V$ 同时成立”来确定的，而不是把矩阵直接看成 $\fg$ 的元素；这一步在特征 $p$ 下会失效。]
这个公式的含义如下：(1.2-4) 右端出现的表达式 $gA_ig^{-1}-(\partial_{z_i}g)g^{-1}$，对 $G$ 的任意有限维表示 $V$ 都是 $\End V$ 中确定的元素。于是上面讨论的 Tannakian 形式表明，存在 $\fg$ 的一个“万有”元素，它在 $V$ 上的作用由这个公式给出。

\subsection{正则奇点与非正则奇点}

与 $GL_n$ 的情形一样，复代数簇 $X$ 上带全纯平坦联络的全纯主 $G$-丛给出从 $X$ 的基本群到 $G$ 的同态的一个等价类。但这是否建立了相应等价类之间的双射？若 $X$ 紧，确实如此；但若 $X$ 不紧，平坦丛就比 $\pi_1(X)$ 的表示多。为了得到双射，需要对联络加一个条件：要求它在无穷远处具有正则奇点。

若 $X$ 是从射影（因而紧）曲线 $\overline{X}$ 挖去有限多个点得到的曲线，这个条件意味着联络算子 $\partial_z+A(z)$ 在每个被挖去的点处有阶数至多为一的极点。（对高维代数簇 $X$，把联络限制到 $X$ 中的所有曲线上，就直截了当地推广了该条件。）

例如考虑 $X=\mathbf{A}^\times=\Spec\C[t,t^{-1}]$，并设 $E$ 的秩为 $1$，即是线丛。可以把 $X$ 嵌入射影曲线 $\PP^1$，于是“无穷远处”有两个点：$t=0$ 与 $t=\infty$。$\mathbf{A}^\times$ 上的任何全纯线丛都可平凡化。$\mathbf{A}^\times$ 上一般的联络算子于是写成
\begin{equation}
    \nabla_{\partial_t}=\partial_t+A(t),\qquad
    A(t)=\sum_{i=N}^{M}A_it^i.
    \tag{1.2-5}
\end{equation}
$\mathbf{A}^\times$ 上可逆函数群由元素 $\alpha t^n$ 组成，其中 $\alpha\in\C^\times$，$n\in\Z$。它通过规范变换 (1.2-4) 作用在这样的算子上，而 $\alpha t^n$ 的作用如下：
\[
    A(t)\longmapsto A(t)-\frac{n}{t}.
\]
$\mathbf{A}^\times$ 上带平坦联络的线丛的等价类集合，与算子 (1.2-5) 的集合模去规范变换群所得的商一一对应；等价地，与 $\C[t,t^{-1}]$ 模去群 $\Z\cdot\frac{1}{t}$ 的加法作用所得的商一一对应。由此可见这个集合非常大。

下面只看带正则奇点的联络。联络 (1.2-5) 在 $t=0$ 处有正则奇点，意味着 $N\ge -1$。要看 $t=\infty$ 处的条件，作变量替换 $u=t^{-1}$。

% [Tutor Note: p.17。$t=\infty$ 处的计算：由 $\partial_t=-u^2\partial_u$ 得 $\nabla_{\partial_u}=\partial_u-u^{-2}A(u^{-1})$。把 $A(t)=\sum_{i=N}^{M}A_it^i$ 代入，得 $\sum_i A_iu^{-i-2}$；要在 $u=0$ 处极点阶数不超过一，须 $-i-2\ge-1$，即 $i\le-1$，故 $M\le-1$。两端条件合起来迫使 $A(t)=a/t$。这里容易把 $t$ 与 $u$ 的极点阶数弄反，务必按此步骤核对。]
于是由 $\partial_t=-u^2\partial_u$ 得到
\[
    \nabla_{\partial_u}=\partial_u-u^{-2}A(u^{-1}).
\]
因此无穷远处的正则奇点条件是 $M\le -1$；要同时满足两个条件，$A(t)$ 必须具有形式 $A(t)=\frac{a}{t}$。

把规范变换也考虑进来，我们断定：$\mathbf{A}^\times$ 上带正则奇点平坦联络的线丛的等价类集合同构于 $\C/\Z$。给定满足 $A(t)=\frac{a}{t}$ 的联络 (1.2-5)，方程
\[
    \left(\partial_t+\frac{a}{t}\right)\Phi(t)=0
\]
的解为 $\Phi_C(t)=C\exp(-at)$，其中 $C\in\C$\transnote{此处原文的 $\Phi_C(t)=C\exp(-at)$ 应为 $\Phi_C(t)=Ct^{-a}=C\exp(-a\log t)$，译文照原书保留。代入方程 $(\partial_t+a/t)\Phi=0$ 即可验证：$Ct^{-a}$ 才是解；而 $C\exp(-at)$ 满足的是 $(\partial_t+a)\Phi=0$，其单值性平凡，与紧接其后的 $\exp(-2\pi ia)$ 矛盾。[来源：原书 p.17]}。这些就是我们的线丛的水平截面。于是该联络沿 $\mathbf{A}^\times$ 中逆时针绕过原点的圈的单值性等于 $\exp(-2\pi ia)$。我们得到 $\pi_1(\mathbf{A}^\times)\simeq\Z$ 的一维表示，它把 $1\in\Z$ 送到 $\exp(-2\pi ia)$。因此指数映射建立了 $\mathbf{A}^\times$ 上带正则奇点联络的等价类集合与 $\mathbf{A}^\times$ 基本群的表示之间的双射。

% [Tutor Note: p.17 原文疑似笔误，译文照原书保留。方程 $(\partial_t+a/t)\Phi=0$ 的解应为 $\Phi_C(t)=Ct^{-a}=C\exp(-a\log t)$，而不是 $C\exp(-at)$：后者满足的是 $(\partial_t+a)\Phi=0$，其单值性平凡，与紧接着断言的单值性 $\exp(-2\pi ia)$ 矛盾。报告已登记。]
但若允许非正则奇点，就会得到多得多的联络。例如考虑联络 $\partial_t+\frac{1}{t^2}$。它在 $t=0$ 处有非正则奇点，对应的水平截面形如 $C\exp(-1/t)$，$C\in\C$；于是该联络对应的 $\pi_1(\mathbf{A}^\times)$ 单值性表示是平凡的，与联络 $\nabla_{\partial_t}=\partial_t$ 相同。然而在代数规范变换的作用下，这个联络并不等价于联络 $\partial_t$。当然可以用函数 $\exp(-1/t)$ 的规范作用把它们互变，但这个函数不是代数的。因此，若允许非正则奇点，代数规范等价类就比基本群到 ${}^L G$ 的同态多得多。

% [Tutor Note: p.17 背景：这个例子说明为什么必须在“代数规范变换”下取商。$\exp(-1/t)$ 是解析可逆函数但不是 $\C((t))$ 中的元素，故它不属于 ${}^L G((t))$；忽略这一点就会错误地认为非正则奇点可由正则奇点经规范变换得到（即所谓 Stokes 现象消失）。在 $\F_q((t))$ 一侧，对应的现象是野惯性。]
\subsection{联络作为 Langlands 参数}

% [Tutor Note: p.18。本小节给出全章的核心定义：$\Loc_{{}^L G}(D^\times)$。它把“同态 $\pi_1(D^\times)\to{}^L G$ 的等价类”替换为“${}^L G$-丛加任意阶极点联络的规范等价类”，从而把离散的参数集合换成无穷维的几何对象。记住这个定义，第 10 章的对应陈述（本单元 p.30 末）将直接引用它。]
我们的目标是给同态 $\pi_1(D^\times)\to{}^L G$ 的等价类集合找一个几何上的增强，因为这样的同态太少。上面的讨论提示了一条可行途径。我们看到，这样的同态等同于 $D^\times$ 上带一个在原点处有正则奇点的联络的 ${}^L G$-丛。这里“去心圆盘”指概形 $D^\times=\Spec\C((t))$，因此按定义 $D^\times$ 上的函数代数就是 $\C((t))$。$D^\times$ 上的任何 ${}^L G$-丛都可平凡化\transnote{$D^\times=\Spec\C((t))$ 上的 ${}^L G$-丛必可平凡化：域上的约化群主丛由 $H^1(\C((t)),{}^LG)$ 分类，而 $\C((t))$ 是 $C_1$ 域（Tsen 定理），由 Steinberg 定理该上同调群平凡。正因如此，“$D^\times$ 上带联络的 ${}^L G$-丛”才等同于“${}^L G((t))$ 中的一阶算子模规范变换”，即式 (1.2-7)。[来源：原书 p.17--18；J.-P. Serre, \emph{Galois Cohomology}, Ch.~III \S2]}，所以在其上定义联络就归结为给出一个一阶算子
\begin{equation}
    \partial_t+A(t),\qquad A(t)\in{}^L\fg((t)),
    \tag{1.2-6}
\end{equation}
其中 ${}^L\fg$ 是 Langlands 对偶群 ${}^L G$ 的李代数。

改变平凡化相当于 (1.2-4) 中 $g\in{}^L G((t))$ 的规范变换。$D^\times$ 上带联络的 ${}^L G$-丛的等价类集合，与算子 (1.2-6) 的规范等价类集合一一对应。把这个集合记作 $\Loc_{{}^L G}(D^\times)$。于是
\begin{equation}
    \Loc_{{}^L G}(D^\times)
    =\left\{\partial_t+A(t),\ A(t)\in{}^L\fg((t))\right\}\big/{}^L G((t)).
    \tag{1.2-7}
\end{equation}

联络 (1.2-6) 有正则奇点，当且仅当 $A(t)$ 在 $t=0$ 处极点阶数至多为一。若 $A_{-1}$ 是 $A(t)$ 在 $t=0$ 处的留数，则相应的绕原点单值性是 ${}^L G$ 中的元素，等于 $\exp(2\pi iA_{-1})$\transnote{此处符号与 p.17 的约定相反，原文如此，译文照录。按 p.17 的计算，$A(t)=a/t$ 对应的单值性是 $\exp(-2\pi ia)$；而按 $\Phi=t^{-A_{-1}}$，留数 $A_{-1}$ 给出的单值性应为 $\exp(-2\pi iA_{-1})$。若坚持 $\exp(2\pi iA_{-1})$，则须约定绕原点顺时针绕行；引用本书时请先统一这一约定。[来源：原书 p.17--18 对照]}。两个带正则奇点的联络（代数地）规范等价，当且仅当它们的单值性彼此共轭。因此带正则奇点联络的等价类集合就是 ${}^L G$ 的共轭类集合，等价地，就是同态 $\Z\to{}^L G$ 的集合。这些就是我们先前见到的天真 Langlands 参数（这里忽略了拓扑基本群与代数基本群的差别）。但现在我们允许原点处具有任意奇点，即正则与非正则奇点都可出现，从而得到多得多的规范等价类。

% [Tutor Note: p.18 与 p.17 的单值性符号不一致，原文如此。p.17 对 $A(t)=a/t$ 给出单值性 $\exp(-2\pi ia)$（逆时针绕行），而这里对留数 $A_{-1}$ 给出 $\exp(2\pi iA_{-1})$，符号相反。按 $\Phi=t^{-A_{-1}}$ 的计算应为 $\exp(-2\pi iA_{-1})$；但也可能作者在此处采用顺时针绕行的约定。读者引用时请统一约定，报告已登记此疑点。]
我们常常把 $\Loc_{{}^L G}(D^\times)$ 的点称为 $D^\times$ 上的局部系，这里取局部系概念的 de Rham 版本，即带有在原点处任意阶极点联络的主丛。

于是我们得到如下提法：复域上的局部 Langlands 参数应当是 $\Loc_{{}^L G}(D^\times)$ 的点，即 $D^\times$ 上平坦 ${}^L G$-丛的等价类；更具体地说，是一阶微分算子的规范等价类 (1.2-7)\transnote{$\Loc_{{}^LG}(D^\times)$ 是 $\C$ 上几何 Langlands 对应的局部一端：整体情形把 $\Bun_G(X)$ 上的 D-模与 $\Loc_{{}^LG}(X)$ 上的拟凝聚层对应起来（Beilinson--Drinfeld），局部情形则把 $\hat\fg_\kappa\text{-}\mathrm{mod}$ 按 $\Loc_{{}^LG}(D^\times)$ 分解。注意这里的“局部系”取 de Rham 版本：不是局部常值层的等价类，而是带任意阶极点联络的 ${}^LG$-丛的规范等价类，因而比 $\pi_1$ 的表示多得多。[来源：原书 p.18、p.30；nLab “geometric Langlands correspondence”，访问 2026-11-23]}。

我们要指出，局部域 $\F_q((t))$ 的 Galois 群结构非常复杂：除了来自 $\operatorname{Gal}(\overline{\F}_q/\F_q)$ 的部分与温和惯性（它们类比于我们在复域看到的单值性）之外，还有野惯性子群，其表示极其复杂。一个由来已久的想法是：在复域，带非正则奇点的平坦联络在某种意义上扮演了“限制到野惯性子群上非平凡的 Galois 群表示”的角色。我们的提法只是在 Langlands 对应的背景下利用了这一想法。

\section{圈群的表示}

\begin{keypoints}
    \item \textbf{核心思想：} 待参数化的对象不是向量空间上的光滑表示，而是“范畴化”的表示，即带圈群作用的范畴；这样才能同时适用于有限域与复域。
    \item \textbf{数学工具：} 同余子群与光滑性；$\ell$-adic 层与 Grothendieck 的 fonctions-faisceaux 词典；可构造层与 D-模之间的 Riemann--Hilbert 对应；Beilinson--Bernstein 局部化。
    \item \textbf{关键结论：} 代数（或可微）表示过于贫乏：不可约光滑代数表示必为一维；应改取仿射 Kac--Moody 代数 $\hat{\fg}_\kappa$ 的光滑模范畴，它在 $\Loc_{{}^L G}(D^\times)$ 上的“纤维”给出局部 Langlands 对应的表述。
    \item \textbf{记号提醒：} $K_N$ 是同余子群；$\hat{\fg}_\kappa$ 是水平 $\kappa$ 的仿射 Kac--Moody 代数；临界水平 $\kappa_c=-\frac{1}{2}\kappa_\fg$。
\end{keypoints}

% [Tutor Note: p.18。全节的问题意识：参数一侧已经由 $\Loc_{{}^L G}(D^\times)$ 给定，接下来要确定“被参数化的对象”。作者刻意不走“光滑表示”的老路，而要把表示范畴化，理由在 1.3.1--1.3.3 逐层展开：向量空间上的代数表示太贫乏，而层与范畴能承载几何。]
解决了 Langlands 参数的问题之后，我们必须决定要参数化的对象是什么。回忆在经典情形，同态 $W'_F\to{}^L G$ 参数化了群 $G(F)$ 的不可约光滑表示，其中 $F=\F_q((t))$。我们先把这个概念翻译到圈群的表示论中。

\subsection{光滑表示}

圈群 $G((t))$ 包含同余子群
\begin{equation}
    K_N=\{g\in G[[t]]\mid g\equiv 1\bmod t^N\},\qquad N\in\Z_+.
    \tag{1.3-1}
\end{equation}
把 $G((t))$ 在复向量空间 $V$ 上的表示称为光滑的，是很自然的：若对任意向量 $v\in V$，存在 $N\in\Z_+$ 使 $K_N\cdot v=v$。若像以前一样用子群 $K_N$（$N\in\Z_+$）作单位元的开邻域基来给 $G((t))$ 赋予拓扑，这个条件就可以解释为连续性条件。

但现在我们的群 $G$ 是复李群（不是有限群），所以 $G((t))$ 是无穷维李群。更精确地说，我们把 $G((t))$ 看作 ind-群，即 ind-概形范畴中的群对象。乍看之下，自然会考虑 $G((t))$ 的代数表示。注意 $G((t))$ 由对应于仿射单根的“parahoric”代数群 $P_i$ 生成。对这些子群，代数表示的概念完全说得通。于是称 $G((t))$ 的一个表示是代数的，若它限制到每个 $P_i$ 上都是代数表示。

然而，这条天真的路径通向如下令人沮丧的事实：$G((t))$ 的不可约光滑表示若还是代数的，则必为一维。为此注意，$G((t))$ 的代数表示给出其李代数 $\fg((t))=\fg\otimes\C((t))$ 的表示，其中 $\fg$ 是 $G$ 的李代数。由 $G((t))$ 的代数表示得到的 $\fg((t))$ 的表示称为可积的。$G((t))$ 的光滑表示给出 $\fg((t))$ 的光滑表示，即对任意向量 $v$，存在 $N\in\Z_+$ 使
\begin{equation}
    \fg\otimes t^N\C[[t]]\cdot v=0.
    \tag{1.3-2}
\end{equation}
可以分解 $\fg=\fg_{\mathrm{ss}}\oplus\mathfrak{r}$，其中 $\mathfrak{r}$ 是中心，$\fg_{\mathrm{ss}}$ 是半单部分。因此 $\fg((t))$ 的不可约表示是 $\fg_{\mathrm{ss}}$ 的不可约表示与 $\mathfrak{r}((t))$ 的不可约表示之张量积\transnote{原文此处印作“$\fg_{\mathrm{ss}}$ 的不可约表示”，按上下文应为“$\fg_{\mathrm{ss}}((t))$ 的不可约表示”，译文照原书保留。理由是李代数的直和分解 $\fg((t))=\fg_{\mathrm{ss}}((t))\oplus\mathfrak r((t))$ 中带形式幂级数的是半单部分，常数部分 $\fg_{\mathrm{ss}}$ 不足以张出 $\fg((t))$ 的表示。[来源：原书 p.19]}；由于 $\mathfrak{r}((t))$ 交换，后者的不可约表示必为一维。现在我们有

% [Tutor Note: p.19。这里写“$\fg_{\mathrm{ss}}$ 的不可约表示”，按上下文应理解为 $\fg_{\mathrm{ss}}((t))$ 的不可约表示：$\fg((t))=\fg_{\mathrm{ss}}((t))\oplus\mathfrak{r}((t))$ 才是李代数直和，而 $\fg_{\mathrm{ss}}$ 本身（常数部分）不够。原文疑漏写 $((t))$，译文照原文保留，报告已登记。]
\begin{lemma}\label{lem:u02-smooth-integrable}
    半单李代数 $\fg$ 的圈代数 $\fg((t))$ 的光滑可积表示必为平凡表示。
\end{lemma}

\begin{proof}
    我们沿用 [Beilinson and Drinfeld 1997] 中 3.7.11(ii) 的论证。设 $V$ 是这样的表示。则存在 $N\in\Z_+$ 使 (1.3-2) 成立。取常值子代数 $\fg\subset\fg((t))$ 的 Cartan 分解 $\fg=\fn_-\oplus\fh\oplus\fn_+$（见下文附录 A.3）。设 $H$ 是 $G$ 中对应于 $\fh\subset\fg$ 的 Cartan 子群。按假设，$V$ 是 $H$ 的代数表示，因此它分解为权空间的直和 $V=\bigoplus V_\chi$，$\chi$ 取遍 $H$ 的整权。于是任意向量 $v\in V$ 可分解为和 $v=\sum_{\chi\in\fh^*}v_\chi$，其中 $v_\chi\in V_\chi$。

    现在注意，存在元素 $h\in H((t))$，使得
    \[
        \fn_+\subset h\left(\fg\otimes t^N\C[[t]]\right)h^{-1}.
    \]
    于是向量 $h\cdot v$ 在 $\fn_+$ 作用下不变。由于 $\fn_+$ 的作用与 $H$ 的整权分次相容，我们得到每个向量 $h\cdot v_\chi$ 也在 $\fn_+$ 下不变。考虑 $V$ 中由 $h\cdot v_\chi$ 生成的 $\fg$-子模。按假设，$\fg$ 在 $V$ 上的作用可积为 $G$ 在 $V$ 上的代数表示。因此 $V$ 是带整优权的有限维不可约 $G$-表示的直和。所以除非 $\chi$ 是整优权，否则 $h\cdot v_\chi=0$，从而 $v_\chi=0$。其次注意，$h^{-1}(\fg\otimes t^N\C[[t]])h$ 含 $\fn_-$。对 $h^{-1}\cdot v$ 施以同样的论证，得到除非 $\chi$ 是反优整权，否则 $v_\chi=0$。

    因此 $v=v_0$，于是 $h\cdot v$ 在 $\fh\oplus\fn_+$ 下不变。但由 $h\cdot v$ 生成的 $\fg$-子模是平凡的。故 $h\cdot v$ 是 $\fg$-不变的，从而 $v$ 是 $h^{-1}\fg h$-不变的。然而李代数 $\fg\otimes t^N\C[[t]]$ 与 $h^{-1}\fg h$（其中 $h$ 取遍 $H((t))$ 中使 $h(\fg\otimes t^N\C[[t]])h^{-1}$ 含 $\fn_+$ 的元素）生成整个李代数 $\fg((t))$。因此 $v$ 是 $\fg((t))$-不变的。于是 $V$ 是 $\fg((t))$ 的平凡表示。
\end{proof}

% [Tutor Note: p.20 证明中省略的两处。其一，为什么存在 $h\in H((t))$ 使 $\fn_+\subset h(\fg\otimes t^N\C[[t]])h^{-1}$：这只要取 $h$ 使 $\operatorname{Ad}_{h^{-1}}$ 把 $\fn_+$ 的根向量压入足够深的 $t$-幂，可用 $H((t))$ 中取值于 $\fh$ 的洛朗级数显式构造。其二，为什么这两个李代数生成整个 $\fg((t))$：因为 $\fg\otimes t^N\C[[t]]$ 与所有共轭 $h^{-1}\fg h$ 一起包含 $\fg\otimes\C[[t]]$ 与 $\fg\otimes t^{-1}\C[t^{-1}]$ 的足够多元素，局部地由 $\fsl_2$-三元组生成。原证明只引 [Beilinson and Drinfeld 1997]，细节需回查该文献。]
由此可见，圈群的代数表示这一类太窄。我们可以放宽条件，考虑可微表示，即把 $G((t))$ 看作李群时的表示。但容易看出结果还是一样。把 $G((t))$ 换成它的中心扩张 $\widehat G$ 也帮不上多少忙：$\widehat G$ 的不可约可积表示由整优权参数化，且它们之间没有扩张 [Kac 1990]。这些表示依然太稀疏，无法用上面考虑的几何数据来参数化。因此我们应当寻找其他类型的表示。

回到局部 Langlands 对应的原始设定。那里考虑的是 $G(\F_q((t)))$ 在 $\C$-向量空间上的表示，所以我们根本用不上 $G(\F_q((t)))$ 作为 $\F_q$ 上 ind-群的代数结构。因此不能指望复圈群 $G((t))$ 的代数（或可微）表示这一类，在 Langlands 对应的观点下有意义。我们应把圈群 $G((t))$ 看成抽象拓扑群，拓扑由同余子群给出；换言之，把它作为抽象群来考虑其光滑表示。

举一个这样的表示为例。考虑由有限线性组合 $\sum\delta_x$ 构成的向量空间，其中 $x$ 取遍商 $G((t))/K_N$ 的点。群 $G((t))$ 自然作用在这个空间上：$g\cdot(\sum\delta_x)=\sum\delta_{g\cdot x}$。对任意 $g\in G((t))$，当 $M$ 足够大时子群 $gK_Ng^{-1}$ 包含 $K_M$，所以这确实是 $G((t))$ 的光滑表示。但它肯定不是代数表示（也不是可微表示）。事实上，它与 $G((t))$ 的代数结构毫无关系，这是严重的缺陷。毕竟，把 Langlands 对应从有限域推广到复域的全部意义，就在于能使用复代数几何的强大工具。若表示与几何完全无关，能从中得到的东西就不多。

% [Tutor Note: p.21。这个 $\sum\delta_x$ 表示其实就是从 $K_N$ 的平凡表示诱导得到的表示（紧诱导），它是光滑的但“完全不几何”。作者用这个反例说明：只靠抽象群的光滑表示论无法把复几何搬进来，必须换成层或范畴。注意这里 $K_N$ 的选取不影响例子的要旨。]
所以我们需要寻找某种几何对象，它把群的表示封装起来，并且在有限域与复域上都有意义。

\subsection{从函数到层}

先重新审视群 $G(F)$（$F=\F_q((t))$）的光滑表示。我们通过矩阵系数把这类表示更具体地实现出来。设 $(V,\pi)$ 是 $G(F)$ 的不可约光滑表示。定义反变（contragredient）表示 $V^\vee$ 为对偶表示 $V^*$ 中所有光滑向量张成的线性包。这个线性包在 $G(F)$ 作用下稳定，因而给出 $G(F)$ 的光滑表示 $(V^\vee,\pi^\vee)$。现在取 $V^\vee$ 中的 $K_N$-不变向量 $\phi$，定义线性映射
\[
    V\longrightarrow C(G(F)/K_N),\qquad v\longmapsto f_v,
\]
其中 $f_v(g)=\langle\pi^\vee(g)\phi,v\rangle$。这里 $C(G(F)/K_N)$ 表示 $G(F)/K_N$ 上 $\C$ 值局部常值函数构成的向量空间。群 $G(F)$ 按公式 $(g\cdot f)(h)=f(g^{-1}h)$ 自然作用在该空间上；上述映射是表示的同态，当 $(V,\pi)$ 不可约时它非零，因而单射。

于是我们把表示实现为商 $G(F)/K_N$ 上函数空间中的表示。更一般地，可以取 $V$ 中 $K$ 的有限维子表示（而非平凡子表示），从而把表示实现为商 $G((t))/K$ 上、取值为某个有限维向量空间的函数空间中的表示。

这里有一个重要的观察：当 $F=\F_q((t))$、$K$ 是 $G(F)$ 的紧子群时，$G(F)/K$ 不只是一个集合，它是定义在 $\F_q$ 上的代数簇（更准确地说，ind-概形）的点集。例如取 $K_0=G(\F_q[[t]])$（极大紧子群），商 $G(F)/K_0$ 就是称为仿射 Grassmann 流形的 ind-概形的 $\F_q$-点集。

接下来回顾一个可以追溯到 Grothendieck 的重要思想：定义在 $\F_q$ 上的代数簇 $X$ 的 $\F_q$-点集上的函数，常常可以看作 $X$ 上所谓 $\ell$-adic 层的“影子”。

% [Tutor Note: p.21--22 背景。Grothendieck 的 fonctions-faisceaux 词典把“层”翻译成“函数”，但这不是范畴等价：不同的层（甚至复形）可以给出同一个函数。正因如此，作者才说层是比函数更“厚”的对象——它保留了函数的来源信息（Frobenius 在每个茎上的作用），从而在复域有对应物。]
下面简单讨论。设 $\ell$ 是不整除 $q$ 的素数。$X$ 上 $\ell$-adic 层范畴的定义分几步（参见 [Milne 1980; Freitag and Kiehl 1988]）。首先考虑 étale 拓扑（其中“开子集”的角色由 étale 态射 $U\to X$ 扮演）下的局部常值 $\Z/\ell^m\Z$-层。$X$ 上的 $\Z_\ell$-层按定义是一族满足自然相容性的局部常值 $\Z/\ell^m\Z$-层 $(\cF_m)$。然后在 $\Z_\ell$-层范畴中消去挠层，得到 $\Q_\ell$-层范畴。类似地，对 $\Q_\ell$ 的有限扩张 $E$，定义 $X$ 上 $E$-层的范畴。最后取 $X$ 上 $E$-层范畴的正向极限，这个范畴的对象称为 $X$ 上局部常值 $\ell$-adic 层。秩为 $n$ 的这样的层，等同于 $X$ 上函数域的 Galois 群的处处无分歧的 $n$ 维表示。因此，局部常值 $\ell$-adic 层是复代数簇情形下通常局部系（相对于解析拓扑）的类比物。

可以把这个 $\ell$-adic 局部系的定义推广：允许 $\Z/\ell^m\Z$-层 $\cF_m$ 是可构造的，即存在 $X$ 由局部闭子簇 $X_i$ 给出的分层，使得层 $\cF|_{X_i}$ 局部常值。这样就得到 $X$ 上可构造 $\ell$-adic 层的概念，简称 $\ell$-adic 层。

从几何上重新表述这一概念的关键一步，是 Grothendieck 的 fonctions-faisceaux 词典（参见 [Laumon 1987]）。设 $\cF$ 是 $\ell$-adic 层，$x$ 是 $X$ 的 $\F_{q_1}$-点，其中 $q_1=q^m$。则有 Frobenius 共轭类 $\mathrm{Fr}_x$ 作用在 $\cF$ 在 $x$ 处的茎 $\cF_x$ 上。于是可以定义 $X$ 的 $\F_{q_1}$-点集上的函数 $f_{q_1}(\cF)$，它在 $x$ 处的值为 $\Tr(\mathrm{Fr}_x,\cF_x)$。这个函数取值于 $\Q_\ell$ 的代数闭包 $\overline{\Q}_\ell$。但取值为 $\overline{\Q}_\ell$ 的函数与取值为 $\C$ 的函数差别不大：两者基数相同，$\overline{\Q}_\ell$ 与 $\C$ 可作为抽象域等同。此外，在多数有趣的情形，取值其实落在 $\overline{\Q}$ 中，而 $\overline{\Q}$ 同时含于 $\overline{\Q}_\ell$ 与 $\C$。

更一般地，若 $\cF$ 是 $\ell$-adic 层的复形，则取 $\mathrm{Fr}_x$ 在 $\cF$ 在 $x$ 处茎上同调的作用的迹的交错和，定义 $X(\F_{q_1})$ 上的函数 $f_{q_1}(\cF)$。映射 $\cF\mapsto f_{q_1}(\cF)$ 把层上的自然运算与函数上的自然运算联系起来（见 [Laumon 1987] 第 1.2 节）。例如，层的拉回对应函数的拉回；带紧支集的层的推前对应函数的纤维积分（这由 Grothendieck--Lefschetz 迹公式得到）。

% [Tutor Note: p.22--23。这里只用“层复形”而不是单个层，是为了让推前等运算保持良好性质；Grothendieck--Lefschetz 迹公式是“带紧支集推前对应纤维积分”的唯一依据，也是 Weil 猜想证明的核心工具。读者若只关心表示论一侧，可先接受这个词典，细节见 [Laumon 1987]。]
设 $K_0(\cS h_X)$ 是 $X$ 上 $\ell$-adic 层范畴的复化 Grothendieck 群。则上述构造给出映射
\[
    K_0(\cS h_X)\longrightarrow \prod_{m\ge 1}X(\F_{q^m}),
\]
并且已知这个映射是单射（见 [Laumon 1987]）。

因此我们可以期望：实现我们表示的那些商 $G(F)/K_N$ 上的函数，会通过这一构造来自 $X$ 上的 $\ell$-adic 层，或更一般地来自 $X$ 上的 $\ell$-adic 层复形。

现在，可构造层的概念（不同于函数的概念）对复代数簇 $X$ 有透明而有意义的类比物：即 $X$ 上那些 $\C$-向量空间层，它们在 $X$ 的某个分层的各层上局部常值。仿射 Grassmann 流形以及作为商 $G(F)/K_N$ 底层的更一般的 ind-概形，在 $\F_q$ 上和 $\C$ 上都有定义。于是自然把这样的层（更精确地说，它们的导出范畴）在 $\C$ 上的这些 ind-概形上的范畴，看作“在其点集上实现 $G(F)$ 的光滑表示的向量空间”的替代物。

% [Tutor Note: p.23 含糊处。$K_0(\cS h_X)\to\prod_m X(\F_{q^m})$ 的单射性（[Laumon 1987]）说明“层复形的类”可以无信息损失地转成“所有有限域扩张上的函数族”；这是把有限域一侧的表示翻译到层一侧的理论依据。注意这里写的是乘积 $\prod$ 而不是不交并：每个 $m$ 给出一个函数族分量。]
于是我们自然来到 [Frenkel and Gaitsgory 2006c] 中提出的想法：我们所需考虑的圈群 $G((t))$ 的表示，不是实现在向量空间上，而是实现在范畴上，例如仿射 Grassmann 流形上凝聚层的导出范畴。当然，这样的范畴有 Grothendieck 群，群 $G((t))$ 也会作用在 Grothendieck 群上，从而给出 $G((t))$ 在向量空间上的表示。但考察范畴表示会得到多得多的结构：范畴的对象以及作用都有几何意义，因此我们可以尽可能多地使用几何。

把上述概括一下：对每个局部 Langlands 参数 $\chi\in\Loc_{{}^L G}(D^\times)$，我们希望配上一个范畴 $\cC_\chi$，并带上圈群 $G((t))$ 的作用。但这些 $\cC_\chi$ 应当是什么样的范畴，我们又期待它们满足什么性质？

为了更接近回答这些问题，我们想再讨论上面讨论中还能走的两步，以到达本书要考虑的“带圈群作用的范畴”的类型。

\subsection{一个玩具模型}

% [Tutor Note: p.23--24 结构。玩具模型把 $G(\F_q((t)))$ 换成 $G(\F_q)$、把圈群表示换成 $G(\F_q)$ 在 $G(\F_q)/B(\F_q)$ 上函数空间中的表示，再用旗簇上的 D-模替代函数。这条链条是全书“几何化”的模板：函数 → 层 → D-模 → 范畴。]
此时稍作绕道、考察我们的构造的一个玩具模型，是有启发性的。设 $G$ 是 $\Z$ 上的分裂约化群，$B$ 是其 Borel 子群。$G(\F_q)$ 的一个自然表示实现在商 $G(\F_q)/B(\F_q)$ 上复值（或 $\overline{\Q}_\ell$ 值）函数的空间中。自然要问：若把域 $\F_q$ 换成复域、把群 $G(\F_q)$ 换成 $G(\C)$，这个表示的“正确”类比物是什么？这可以看作我们困境的简化版，因为现在看的是 $G(\F_q)$ 而不是 $G(\F_q((t)))$。

商 $G(\F_q)/B(\F_q)$ 是定义在 $\Z$ 上的代数簇的 $\F_q$-点集，这个代数簇称为 $G$ 的旗簇，记作 $\Fl$。上一节的讨论提示：首先要把旗簇 $\Fl(\F_q)$ 上函数的概念，换成簇 $\Fl_{\F_q}=\Fl\otimes_\Z\F_q$ 上 $\ell$-adic 层的概念。

其次，把 $\Fl$（看作 $\F_q$ 上代数簇）上 $\ell$-adic 层的概念，换成 $\Fl_\C=\Fl\otimes_\Z\C$（$\C$ 上代数簇）上可构造层的概念。复代数群 $G_\C$ 自然作用在 $\Fl_\C$ 上，从而作用在这个范畴上。现在再对这个范畴作两种重新表述。

首先，对光滑复代数簇 $X$，有 Riemann--Hilbert 对应：它是 $X$ 上可构造层的导出范畴与 $X$ 上全纯且具正则奇点的 holonomic D-模的导出范畴之间的等价。

这里我们考虑 $X$ 上代数微分算子层及其上的模层，简称 D-模。D-模最简单的例子，是 $X$ 上某个带平坦联络的向量丛的截面层。平坦联络使我们能用函数乘任何截面，也能用平坦联络让向量场作用于截面。这两个作用生成微分算子层在丛的截面上的作用。于是这个丛的水平截面层就是 $X$ 上的局部常值层。上面已经看到，$X$ 上秩为 $n$、联络有正则奇点的丛的同构类，与 $X$ 上秩为 $n$ 的局部常值层的同构类（等价地，$\pi_1(X)$ 的 $n$ 维表示）之间存在双射。这个双射可以提升为相应范畴的等价，而一般的 Riemann--Hilbert 对应就是这一范畴等价的推广，涵盖更一般的 D-模。

Riemann--Hilbert 对应使我们能给 $X$ 上任何 holonomic D-模配上 $X$ 上可构造层的复形，由此得到相应导出范畴之间的函子；若只限于具正则奇点的 holonomic D-模，这个函子是等价（详见 [Borel et al. 1987; Gelfand and Manin 1994]）\transnote{Riemann--Hilbert 对应是“局部系 $\leftrightarrow$ 平坦联络”这古典对应的范畴化：左边是可构造层（沿某个分层局部常值）的导出范畴，右边是正则 holonomic D-模的导出范畴。其中“正则”不可省：去掉它只能得到形式 D-模，单值性不足以刻画对象，还需补上 Stokes 结构。[来源：原书 p.24 及所引 Borel et al. 1987、Gelfand--Manin 1994]}。

于是，在 $\C$ 上可以从可构造层过渡到 D-模。在我们的情形，考虑旗簇 $\Fl_\C$ 上（正则 holonomic）D-模的范畴。这个范畴带有 $G_\C$ 的自然作用。

% [Tutor Note: p.24。Riemann--Hilbert 对应有两个必须记住的限制：holonomic 与正则奇点。去掉正则奇点，可构造层只能对应到“形式地”的 D-模，单值性不足以完全刻画（对应的是 Stokes 结构）；去掉 holonomic，则只能得到导出范畴之间的函子而非等价。后文 1.3.4 使用 D-模时都默认正则 holonomic 的情形。]
最后注意，$G_\C$ 的李代数 $\fg$ 通过向量场无穷小地作用在旗簇上。因此给定 $\Fl_\C$ 上的 D-模 $\cF$，其整体截面空间 $\Gamma(\Fl_\C,\cF)$ 具有 $\fg$-模结构。这样得到从 $\Fl_\C$ 上 D-模范畴到 $\fg$-模范畴的函子 $\Gamma$。A. Beilinson 与 J. Bernstein 证明了：这个函子是从 $\Fl_\C$ 上一切 D-模（不必正则 holonomic）的范畴，到 $\fg$-模范畴 $\cC_0$ 的等价；这里 $\cC_0$ 由那些使泛包络代数 $U(\fg)$ 的中心通过增广特征标作用的 $\fg$-模组成\transnote{这就是 Beilinson--Bernstein 局部化：旗簇上的 D-模范畴等价于“中心按增广特征标作用”的 $\fg$-模范畴 $\cC_0$，而整体截面函子 $\Gamma$ 给出这个等价。它的意义在于把“群在函数空间上的表示”换成“范畴连同群作用”，正是本节要的几何类比物；$\cC_0$ 中的对象（如 Verma 模）因此获得几何解释。[来源：A. Beilinson, J. Bernstein, \emph{C. R. Acad. Sci. Paris} \textbf{292} (1981), 15--18]}。

% [Tutor Note: p.25 背景。Beilinson--Bernstein 局部化把“旗簇上的 D-模”与“具平凡中心特征的 $\fg$-模”等同起来，它是几何表示论的出发点。$\cC_0$ 中平凡中心特征意味着 $U(\fg)$ 的中心按增广特征标（即中心元素的常值项）作用；在这一情形，Verma 模的中心特征由最高权的“整性”决定。这个等价是“范畴化表示”的第一个具体范例。]
于是，现在可以回答我们的问题：群 $G(\F_q)$ 在商 $G(\F_q)/B(\F_q)$ 上函数空间中的表示，其有意义的几何类比物是什么？答案如下：它是一个范畴，带上代数群 $G_\C$ 的作用。这个范畴有两种化身：一是旗簇 $\Fl_\C$ 上 D-模的范畴，二是 $\fg$ 的具平凡中心特征的模范畴 $\cC_0$。两个范畴都带有群 $G_\C$ 的自然作用。

下面暂停一下，把“群 $G_\C$ 作用于范畴 $\cC_0$”的确切含义说清楚。为简单起见，我们描述相应的 $\C$-点群 $G(\C)$ 的作用。\footnote{更一般地，对任意 $\C$-代数 $R$，都有 $G(R)$ 作用在相应的、在 $R$ 上基变换后的范畴上。于是我们自然被引向“代数群（或更一般地，群概形）作用在阿贝尔范畴上”的概念，其精确表述见 [Frenkel and Gaitsgory 2006c] 第 20 节。} 意思是：每个元素 $g\in G$ 给出 $\cC_0$ 上的函子 $F_g$，使得 $F_1$ 是恒等函子，且函子 $F_{g^{-1}}$ 是 $F_g$ 的拟逆。此外，对任意一对 $g,h\in G$，有固定的函子同构 $i_{g,h}:F_{gh}\to F_g\circ F_h$，使得对任意三元组 $g,h,k\in G$，同构 $F_{ghk}\to F_g\circ F_h\circ F_k$ 满足等式 $i_{h,k}i_{g,hk}=i_{g,h}i_{gh,k}$。（我们指出，最后一个条件可以放松：只要求 $i_{h,k}i_{g,hk}=\gamma_{g,h,k}i_{g,h}i_{gh,k}$，其中对每个三元组 $g,h,k\in G$，$\gamma_{g,h,k}$ 是非零复数；这些数随后须满足三上循环条件。但下文只考虑 $\gamma_{g,h,k}\equiv 1$ 的情形。）\transnote{这是“群作用在范畴上”的严格定义：不仅每个 $g$ 要给出函子 $F_g$，还要求同构 $i_{g,h}:F_{gh}\to F_g\circ F_h$ 满足三上循环相容性，否则只是“相差一个上循环”的射影作用。把 $i_{h,k}i_{g,hk}=\gamma_{g,h,k}i_{g,h}i_{gh,k}$ 中的 $\gamma$ 放松为非零常数，就得到射影作用；本书只处理 $\gamma\equiv1$ 的严格情形，这样后面才能把 $\hat\fg_\kappa\text{-}\mathrm{mod}$ 上的作用当作真正的群作用。[来源：原书 p.26 脚注 2 所引 Frenkel--Gaitsgory 2006c 第 20 节]}

函子 $F_g$ 的定义如下。给定 $\fg$ 的表示 $(V,\pi)$ 与元素 $g\in G(\C)$，定义新表示 $F_g((V,\pi))=(V,\pi_g)$，其中按定义 $\pi_g(x)=\pi(\operatorname{Ad}_g(x))$。设 $(V,\pi)$ 不可约。则容易看出 $(V,\pi_g)\simeq(V,\pi)$ 当且仅当 $(V,\pi)$ 可积，即由 $G$ 的代数表示得到；这等价于该表示有限维。但一般表示 $(V,\pi)$ 是无穷维的，所以它不会同构于 $(V,\pi_g)$，至少对某些 $g\in G$ 如此。

现在考虑 $\cC_0$ 中的态射，它们就是 $\fg$-同态。给定表示 $(V,\pi)$ 与 $(V',\pi')$ 之间的 $\fg$-同态，即线性映射 $T:V\to V'$ 使对一切 $x\in\fg$ 有 $T\pi(x)=\pi'(x)T$，我们令 $F_g(T)=T$。此时同构 $i_{g,h}$ 全等于恒等。

% [Tutor Note: p.26 易错处。$G$ 作用于范畴的数据不只是“每个 $g$ 给一个函子”，还要求 $i_{g,h}$ 满足三上循环相容性；若只要求相差一个非零常数 $\gamma_{g,h,k}$，则得到“射影”作用（$\gamma$ 是三上循环）。本书只考虑 $\gamma\equiv1$ 的严格作用。另外，$F_g$ 把 $(V,\pi)$ 送到 $(V,\pi\circ\operatorname{Ad}_g)$ 的做法，就是把“共轭表示”视为群作用；这正是 1.3.3 之后对 $\hat{\fg}_\kappa$-mod 作同样构造的模板。]
范畴 $\cC_0$ 中最简单的对象，是从 Borel 子代数 $\fb\subset\fg$ 的一维表示诱导得到的 Verma 模。相应的 D-模是“$\delta$-函数”型 D-模，支撑在 $\Fl_\C$ 中被 $\fb$ 稳定的点上。在范畴 $\cC_0$ 的 Grothendieck 群中，这些对象的类张成 $G(\C)$（此时看作离散群！）的一个子表示，它看起来恰好就像我们在 1.3.2 节开头定义的那个表示。我们做到的是：把这个表示换成了从复代数群 $G$（而不是相应的离散群）的表示论观点看有意义的东西，即范畴 $\cC_0$。

\subsection{回到圈群}

% [Tutor Note: p.26。这里出现第二个“引理 1.3.1”（原文编号重复）：按所在节推断疑应为 1.3.2，而 p.28 处印作 1.3.2 的引理按序疑应为 1.3.3。译文按环境自动编号，不手写编号；编号问题已登记在报告。]
在我们为局部 Langlands 对应寻找复类比物的征途中，必须决定用什么替代群 $G(F)$（$F=\F_q((t))$）的光滑表示这一概念。如上一节的讨论所示，应当考虑复圈群 $G((t))$ 在 ind-概形 $G((t))/K$ 上各种 D-模范畴中的表示，其中 $K$ 是 $G((t))$ 的“紧”子群，如 $G[[t]]$，或 Iwahori 子群（同态 $G[[t]]\to G$ 下 Borel 子群 $B\subset G$ 的原像）；或者考虑李代数 $\fg((t))$ 的表示范畴。两种方案都可行，都会引出有趣的结果与猜想，我们将在第 10 章按 [Frenkel and Gaitsgory 2006c] 详细讨论。本书将集中于第二种方案，考虑圈代数 $\fg((t))$ 上的（投射）模范畴。

群 $G((t))$ 按上一节描述的方式作用在 $\fg((t))$ 的表示范畴上。为使 $G((t))$ 的相应表示光滑，需要把注意力限制在那些使李子代数 $\fg\otimes t^N\C[[t]]$ 的作用对某个 $N>0$ 可积的表示上。确实，在这样的表示上，$\fg\otimes t^N\C[[t]]$ 的作用可以指数化为其李群即同余子群 $K_N$ 的作用。若 $(V,\pi)$ 是这样的表示，则对每个 $g\in K_N$，$g$ 在 $V$ 上的作用算子给出 $(V,\pi)$ 与 $F_g((V,\pi))$ 之间的同构。因此可以说 $(V,\pi)$ 在 $F_g$ 下“稳定”。

如下引理表明，这个条件本质上等价于 $\fg((t))$-模光滑，即任何向量都被某个足够大的 $M$ 对应的李代数 $\fg\otimes t^M\C[[t]]$ 零化。

\begin{lemma}\label{lem:u02-smooth-fg}
    设 $(V,\pi)$ 是李代数 $\fg((t))$ 上的光滑有限生成模。则存在 $N>0$，使李子代数 $\fg\otimes t^N\C[[t]]$ 在 $V$ 上的作用可积。
\end{lemma}

\begin{proof}
    设 $v_1,\dots,v_k$ 是 $V$ 的一组生成元。由于 $(V,\pi)$ 光滑，每个 $v_i$ 被某个李子代数 $\fg\otimes t^{N_i}\C[[t]]$ 零化。令 $M_i$ 是诱导表示
    \[
        M_i=\Ind_{\fg\otimes t^{N_i}\C[[t]]}^{\fg((t))}\C
        =U(\fg((t)))\otimes_{\fg\otimes t^{N_i}\C[[t]]}\C.
    \]
    则有满同态 $\bigoplus_{i=1}^{k}M_i\to V$，把 $M_i$ 的生成元送到 $v_i$，$i=1,\dots,k$。取 $N$ 为 $N_1,\dots,N_k$ 中的最大数。则 $\fg\otimes t^N\C[[t]]$ 在每个 $M_i$ 上的作用可积，从而在 $V$ 上也如此。
\end{proof}

% [Tutor Note: p.27。证明的关键是：光滑性只保证每个生成元各被某个 $N_i$ 零化，有限生成性则把这些 $N_i$ 的最大者 $N$ 变成公共界；诱导模 $M_i$ 把“在一个子代数上平凡”翻译成“可积”，于是满同态的像 $V$ 也可积。注意这里只需 $V$ 有限生成，无须可数生成，故最大数存在。]

\subsection{从圈代数到其中心扩张}

于是，可以把“局部非阿基米德域上约化群的光滑表示”的范畴类比物，取为：圈代数 $\fg((t))$ 上光滑有限生成模的范畴，并带上圈群 $G((t))$ 的自然作用。（下文中我们将去掉有限生成这个条件。）


不过要注意，也可以改取 $\fg((t))$ 的某个中心扩张的光滑表示范畴。群 $G((t))$ 仍通过伴随作用作用在这样的中心扩张上。由于 $G((t))$ 在范畴上的作用是通过其伴随作用得到的，用中心子代数扩张 $\fg((t))$ 不会造成损害。

李代数的中心扩张概念在附录 A.4 中有详细描述。特别地，那里说明：$\fg((t))$ 的中心扩张的等价类由第二上同调 $H^2(\fg((t)),\C)$ 描述。\footnote{这里须考虑我们在圈代数 $\fg((t))$ 上取的拓扑：零元的开邻域基由李子代数 $\fg\otimes t^N\C[[t]]$（$N\in\Z_+$）给出。因此应限于相应的连续上同调。下文对这个连续上同调仍使用同一记号 $H^2(\fg((t)),\C)$。}

我们用分解 $\fg=\fg_{\mathrm{ss}}\oplus\mathfrak{r}$，其中 $\mathfrak{r}$ 是中心，$\fg_{\mathrm{ss}}$ 是半单部分。则可以证明
\[
    H^2(\fg((t)),\C)\simeq H^2(\fg_{\mathrm{ss}}((t)),\C)\oplus H^2(\mathfrak{r}((t)),\C).
\]
换句话说，$\fg((t))$ 的中心扩张由其到 $\fg_{\mathrm{ss}}((t))$ 与到 $\mathfrak{r}((t))$ 的限制决定。我们将考虑的中心扩张在交换部分 $\mathfrak{r}((t))$ 上是平凡的，所以我们的中心扩张将直和分解为 $\hat{\fg}_{\mathrm{ss}}\oplus\mathfrak{r}((t))$。这个李代数的不可约表示同构于 $\hat{\fg}_{\mathrm{ss}}$ 的不可约表示与 $\mathfrak{r}((t))$ 的不可约表示之张量积。但李代数 $\mathfrak{r}((t))$ 交换，其不可约表示一维。这些表示不难单独处理，所以从今以后我们集中于 $\fg((t))$ 的中心扩张的表示，其中 $\fg$ 是半单李代数。

% [Tutor Note: p.27--28 易错处。$H^2(\fg((t)),\C)$ 的分解与“中心扩张限制到两个直和项”是一回事，但连续上同调必须用 $\fg\otimes t^N\C[[t]]$ 的拓扑，否则 $H^2$ 会不同。下文 (1.3-3) 的上循环只有一个独立参数，正与此处的“空间一维”相符。]
若 $\fg$ 半单，则可分解为单李代数 $\fg_i$（$i=1,\dots,m$）的直和，并且
\[
    H^2(\fg((t)),\C)\simeq\bigoplus_{i=1}^{m}H^2(\fg_i((t)),\C).
\]
同样，不妨各单因子 $\fg_i$ 分别处理。因此从今以后假定 $\fg$ 是单李代数。此时对 $H^2(\fg((t)),\C)$ 有如下描述。

回忆：李代数 $\fg$ 上的内积 $\kappa$ 称为不变的，若
\[
    \kappa([x,y],z)+\kappa(y,[x,z])=0,\qquad \forall x,y,z\in\fg.
\]
这个公式是公式 $\kappa(\operatorname{Ad}_g\cdot y,\operatorname{Ad}_g\cdot z)=\kappa(y,z)$ 的无穷小版本，后者表示该型在相应李群（$g$ 是李群中的元素）作用下的不变性。有限维单李代数 $\fg$ 上非退化不变内积的向量空间是一维的。从 $\fg$ 的任一非平凡有限维表示 $\rho_V:\fg\to\End V$ 出发，可用公式
\[
    \kappa_V(x,y)=\Tr_V(\rho_V(x)\rho_V(y))
\]
造出这样的型。标准取法是伴随表示，它给出 Killing 型 $\kappa_\fg$。下述结果是众所周知的\transnote{对单李代数 $\fg$，$H^2(\fg((t)),\C)\cong\C$ 是 Gabber--Kac 的定理，它说明圈代数的中心扩张在同构意义下只由一个参数（水平 $\kappa$）决定。这里的 $H^2$ 指连续上同调：拓扑取子代数 $\fg\otimes t^N\C[[t]]$（$N\in\Z_+$）为零点的开邻域基（见脚注 3），若改用普通李代数上同调则会得到大得多的空间。[来源：O. Gabber, V. G. Kac, \emph{Bull. Amer. Math. Soc. (N.S.)} \textbf{5} (1981), 185--189；原书 p.27 脚注 3]}。

\begin{lemma}\label{lem:u02-h2}
    对单李代数 $\fg$，空间 $H^2(\fg((t)),\C)$ 是一维的，并可等同于 $\fg$ 上不变双线性型的空间。给定这样的型 $\kappa$，相应的中心扩张可由上循环
    \begin{equation}
        c(A\otimes f(t),B\otimes g(t))=-\kappa(A,B)\Res_{t=0}f\,dg
        \tag{1.3-3}
    \end{equation}
    构造。
\end{lemma}

这里对形式洛朗幂级数 $a(t)=\sum_{n\in\Z}a_nt^n$，我们约定
\[
    \Res_{t=0}a(t)\,dt=a_{-1}.
\]

% [Tutor Note: p.28 易错处。(1.3-3) 的负号与 $\Res f\,dg$（而不是 $\Res g\,df$）是本书的约定；因 $\Res(f\,dg)=-\Res(g\,df)$（对 $t^{-1}$ 项分部积分），这两点合起来保证 $c$ 反对称。若把负号丢掉，则 (1.3-4) 的中心项符号随之翻转，后面所有水平的符号都会错。本书取 $\kappa_c=-\frac{1}{2}\kappa_\fg$，与这里的符号相配。]
\subsection{仿射 Kac--Moody 代数及其表示}

对应于非零上循环 $\kappa$ 的中心扩张
\[
    0\to\C\mathbf{1}\to\hat{\fg}_\kappa\to\fg((t))\to 0
\]
称为仿射 Kac--Moody 代数。我们把它记作 $\hat{\fg}_\kappa$，并习惯上称 $\kappa$ 为水平。作为向量空间，它等于直和 $\fg((t))\oplus\C\mathbf{1}$，而交换关系为
\begin{equation}
    [A\otimes f(t),B\otimes g(t)]
    =[A,B]\otimes f(t)g(t)-(\kappa(A,B)\Res f\,dg)\mathbf{1},
    \tag{1.3-4}
\end{equation}
其中 $\mathbf{1}$ 是中心元，与其余一切元素交换。注意对非零内积 $\kappa,\kappa'$，李代数 $\hat{\fg}_\kappa$ 与 $\hat{\fg}_{\kappa'}$ 同构。事实上此时 $\kappa=\lambda\kappa'$，其中 $\lambda\in\C^\times$；而映射 $\hat{\fg}_\kappa\to\hat{\fg}_{\kappa'}$（在 $\fg((t))$ 上等于恒等、把 $\mathbf{1}$ 送到 $\lambda\mathbf{1}$）是同构。由引理 \ref{lem:u02-h2}，非零 $\kappa$ 对应的李代数 $\hat{\fg}_\kappa$ 实际上是 $\fg((t))$ 的泛中心扩张。（按定义，李代数 $\hat{\fg}_0$ 是分裂扩张 $\fg((t))\oplus\C\mathbf{1}$。）

% [Tutor Note: p.29 易错处。$\hat{\fg}_\kappa\cong\hat{\fg}_{\kappa'}$ 的同构把 $\mathbf{1}$ 送到 $\lambda\mathbf{1}$，且要求 $\lambda\kappa'=\kappa$；注意方向：$\kappa=\lambda\kappa'$ 时映射从 $\hat{\fg}_\kappa$ 到 $\hat{\fg}_{\kappa'}$。有些文献写法相反，核对时要看 $\lambda$ 乘在哪一边。另外 $\mathbf{1}$ 在不同水平下被“重标度”，这正是后面说“该范畴只依赖 $\kappa$ 的等价类”的原因。]
注意上循环 (1.3-3) 限制到李子代数
\[
    \fg\otimes t^N\C[[t]],\qquad N\in\Z_+,
\]
上等于 $0$，所以它仍是 $\hat{\fg}_\kappa$ 的李子代数。$\hat{\fg}_\kappa$ 的光滑表示，是指使得每个向量都被这个李子代数（对足够大的 $N$）零化的表示。注意把 $\fg((t))$ 换成 $\hat{\fg}_\kappa$ 后，引理 \ref{lem:u02-smooth-integrable} 的陈述仍然成立。

于是定义范畴 $\hat{\fg}_\kappa\text{-}\mathrm{mod}$，其对象是光滑 $\hat{\fg}_\kappa$-模，且中心元 $\mathbf{1}$ 作用为恒等；态射是 $\hat{\fg}_\kappa$ 表示的同态。在本书中，除另有说明，“$\hat{\fg}_\kappa$-模”总是指 $\mathbf{1}$ 作用为恒等的模。\footnote{也可以让 $\mathbf{1}$ 作用为 $\lambda$ 乘恒等，其中 $\lambda\in\C^\times$；但相应的范畴只是与范畴 $\hat{\fg}_{\lambda\kappa}\text{-}\mathrm{mod}$ 等价。}

群 $G((t))$ 对任意 $\kappa$ 都作用在李代数 $\hat{\fg}_\kappa$ 上。确实，$G((t))$ 的中心扩张的伴随作用通过 $G((t))$ 的作用分解。容易算出这个作用，得到
\[
    g\cdot(A(t)+c\mathbf{1})
    =\Bigl(gA(t)g^{-1}+\Res_{t=0}\kappa\bigl((\partial_tg)g^{-1},A(t)\bigr)\mathbf{1}\Bigr).
\]

有意思的是，$\hat{\fg}_\kappa$ 的对偶空间（更确切地说，对偶空间中的一个超平面）可以等同于 $D^\times$ 上平凡 $G$-丛上联络的空间，从而余伴随作用被等同于 $G((t))$ 在这样的联络空间上的规范作用（详见 [Frenkel and Ben-Zvi 2004] 第 16.4 节）。

% [Tutor Note: p.29 背景。这条“余伴随轨道 = 联络空间”的对应是本书后面把中心特征的谱与 $\Loc_{{}^L G}(D^\times)$ 联系起来的起点：对偶空间中的超平面给出水平，而规范作用正是 (1.2-4) 的无穷小版本。注意这里出现的是 $G$ 而不是 ${}^L G$，因为讨论的是中心扩张本身的余伴随几何。]
我们用 $G((t))$ 在 $\hat{\fg}_\kappa$ 上的作用，按 1.3.3 节同样的方式构造 $G((t))$ 在范畴 $\hat{\fg}_\kappa\text{-}\mathrm{mod}$ 上的作用。具体说，设给定对象 $(M,\rho)$，其中 $M$ 是向量空间，$\rho:\hat{\fg}_\kappa\to\End M$ 是使 $M$ 成为光滑 $\hat{\fg}_\kappa$-模的李代数同态。则对每个 $g\in G((t))$，定义 $\hat{\fg}_\kappa\text{-}\mathrm{mod}$ 的新对象 $(M,\rho_g)$，其中按定义 $\rho_g(x)=\pi(\operatorname{Ad}_g(x))$。若有 $\hat{\fg}_\kappa\text{-}\mathrm{mod}$ 中的态射 $(M,\rho)\to(M',\rho')$，则得到显然的态射 $(M,\rho_g)\to(M',\rho'_g)$。于是定义了函子 $F_g:\hat{\fg}_\kappa\text{-}\mathrm{mod}\to\hat{\fg}_\kappa\text{-}\mathrm{mod}$。这些函子满足 1.3.3 节的条件，于是得到 $G((t))$ 在范畴 $\hat{\fg}_\kappa\text{-}\mathrm{mod}$ 上的作用。

% [Tutor Note: p.29 原文笔误，译文照原书保留。此处定义 $\rho_g$ 时原文印作 $\rho_g(x)=\pi(\operatorname{Ad}_g(x))$，但本段并无 $\pi$，对象是 $(M,\rho)$，故应为 $\rho_g(x)=\rho(\operatorname{Ad}_g(x))$。对照 p.25 同一构造写作 $\pi_g(x)=\pi(\operatorname{Ad}_g(x))$，可确认这是排版笔误。报告已登记。]
现在回到 1.2.6 节定义的 Langlands 参数空间 $\Loc_{{}^L G}(D^\times)$。$\Loc_{{}^L G}(D^\times)$ 的元素可以具体描述为：一阶算子 $\partial_t+A(t)$（$A(t)\in{}^L\fg((t))$）模去 ${}^L G((t))$ 作用的规范等价类（见式 (1.2-7)）。

于是可以把复域上的局部 Langlands 对应表述为如下问题：

\begin{quote}
\itshape
对每个局部 Langlands 参数 $\chi\in\Loc_{{}^L G}(D^\times)$，配上一个子范畴 $\hat{\fg}_\kappa\text{-}\mathrm{mod}_\chi\subset\hat{\fg}_\kappa\text{-}\mathrm{mod}$，它在圈群 $G((t))$ 的作用下稳定。
\end{quote}

我们希望把范畴 $\hat{\fg}_\kappa\text{-}\mathrm{mod}$ 看作“纤维化”于局部 Langlands 参数空间 $\Loc_{{}^L G}(D^\times)$ 之上，其中各范畴 $\hat{\fg}_\kappa\text{-}\mathrm{mod}_\chi$ 是“纤维”，群 $G((t))$ 沿这些纤维作用。从这个观点看，范畴 $\hat{\fg}_\kappa\text{-}\mathrm{mod}_\chi$ 应当给出范畴 $\hat{\fg}_\kappa\text{-}\mathrm{mod}$ 在 $\Loc_{{}^L G}(D^\times)$ 上的“谱分解”。

% [Tutor Note: p.30。这里的“谱分解”是比喻性的：只有当 $\hat{\fg}_\kappa\text{-}\mathrm{mod}$ 的中心足够大、其谱覆盖 $\Loc$ 时，才能把范畴按中心特征的纤维分解。临界水平正是这样的情形（下一段）。这也是 Feigin--Frenkel 中心理论（第 2 章起）要建立的结论。]
在本书第 10 章，我们将给出 [Frenkel and Gaitsgory 2006c] 中的具体提法，描述 $\kappa=\kappa_c$ 为临界水平这一特殊情形下的这些范畴；临界水平是 $-\frac{1}{2}\kappa_\fg$，其中 $\kappa_\fg$ 是上面定义的 Killing 型。这个提法基于如下事实：在临界水平，范畴 $\hat{\fg}_{\kappa_c}\text{-}\mathrm{mod}$ 的中心（它等于 $\hat{\fg}_{\kappa_c}$ 的完备泛包络代数的中心）非常大。事实证明，这个中心的谱与 Langlands 参数空间 $\Loc_{{}^L G}(D^\times)$ 密切相关；这将使我们能把 $\hat{\fg}_{\kappa_c}\text{-}\mathrm{mod}_\chi$ 大致定义为具有固定中心特征的光滑 $\hat{\fg}_{\kappa_c}$-模的范畴\transnote{“中心很大”有一条精确的定理：Feigin 与 Frenkel 证明了 Drinfeld 的猜想——临界水平下 $U(\hat\fg_{\kappa_c})$ 的某个完备化的中心同构于 Langlands 对偶李代数 ${}^L\fg$ 所对应的 Gelfand--Dikii 代数（即量子 Drinfeld--Sokolov 约化所得 $W$-代数的极限）。这个中心的谱之所以与 $\Loc_{{}^LG}(D^\times)$ 密切相关，是因为它可以等同于 ${}^LG$ 在 $D^\times$ 上的 oper 空间，而 oper 忘掉约化后正是 $\Loc$ 中的点。[译者按（无外部依据）：oper 与中心谱的等同见本书第 10 章，本单元未展开][来源：B. Feigin, E. Frenkel, \emph{Int. J. Mod. Phys. A} \textbf{7}, Suppl. 1A (1992), 197--215，摘要，访问 2026-11-23]}。

为了更精确地解释这一机制，我们需要发展仿射 Kac--Moody 代数的表示论，特别是描述 $\hat{\fg}_{\kappa_c}$ 的完备泛包络代数中心的结构。下一章我们将开始通往这一目标的漫长旅程；这将占据本书的主要部分。然后第 10 章将按 [Frenkel and Gaitsgory 2004; Frenkel and Gaitsgory 2006c; Frenkel and Gaitsgory 2006a; Frenkel and Gaitsgory 2005; Frenkel and Gaitsgory 2006b; Frenkel and Gaitsgory 2007c] 诸文，说明如何利用这些结果构造圈群的局部几何 Langlands 对应。

% [Tutor Note: p.30 收束。全节把“参数”（1.2 的 $\Loc_{{}^L G}(D^\times)$）与“被参数化对象”（$\hat{\fg}_\kappa\text{-}\mathrm{mod}$ 的子范畴）两端都交代清楚，真正的定理留到第 10 章。读者应记住两条主线：一是临界水平 $\kappa_c=-\frac{1}{2}\kappa_\fg$，二是中心特征的谱与 $\Loc$ 的对应；这是全书后续内容的纲领。]
